\begin{helper}[Base and missing-set control for stable-change paths]
Let $G$ be a finite loopless multigraph with exact edge identities, palette
$[k]$, and distinguished edge $e_0$.  Fix an ETT $\mathcal E$ on $T$
generated by $\varphi$, a \noderef{HierarchyBase} $B$, a
\noderef{HierarchyCounterexample} $C$ on $B$, and a
\noderef{HierarchyMinimumCounterexampleContext} for these data.
Write $H=C.\mathit{hierarchy}$, $q=H.\mathit{dividerCount}$,
$n=\mathcal E.\mathit{rounds}$, $T_n$ for its terminal ladder tree,
$D_n$ for its computed terminal stored set, $R=B.\mathit{closure}$,
$A=\noderef{HierarchyStarredStage}(H,q)$, and
\[
 Y=T|_{\min\{|A|+(|T|-|A|-1),|T|\}}.
\]
Lengths count displayed edges and subtraction is truncated.  Assume
$Y$ elementary under $\varphi$.  Let distinct colors $\alpha,\beta$ be
missing on $Y$, let $Q$ be the full component through a node $z$, with a
\noderef{TwoColorPathComponent} certificate, and let
$\sigma=\noderef{TwoColorComponentSwap}(G,\varphi,\alpha,\beta,z)$.
Put $a=T.\mathit{vertex}(\noderef{HierarchyFirstMissingPosition}(\varphi,T,\alpha))$.
Assume $Q$ disjoint from the full $\alpha,\beta$ component through $a$, and
assume
\[
 (q>0\ \land\ \alpha\in\overline\varphi(H.\mathit{stage}(q)))
 \quad\text{or}\quad
 \bigl(q=0\ \land\
   [\,\alpha\in\overline\varphi(T_n)\ \text{or}\
      (\alpha\in\overline\varphi(A),\ \alpha,\beta\notin D_n)\,]\bigr).
\]
Then $\sigma$ is $(T_n,D_n,\varphi)$-stable if the terminal kind is not PE,
and is \noderef{CombinedStableColoring} on $(T_n,R,D_n,\varphi)$ if it is PE.
Every vertex of $A$ has exactly its old missing-color set.  Moreover
$\sigma$ is proper and leaves exactly $e_0$ uncolored.  Finally, for each
$0\le r\le q$ and $c\in\{\alpha,\beta\}$, if $c$ is both missing on
$A_r=\noderef{HierarchyStarredStage}(H,r)$ and singleton-closed there,
the change is an \noderef{OutsideKempeRecoloring} of $A_r$.
\end{helper}
\begin{proof}
The context supplies properness, one-uncoloredness, all six lower-rung
assertions, fewer-divider induction, lexicographic minimality, and the
large missing-set bound, by
\noderef{HierarchyMinimumCounterexampleContext}.  The counterexample supplies
maximum property, positivity, and the good witness, by
\noderef{HierarchyCounterexample}.  Consecutive exact initial segments
compose by composing their length inequalities and entry equalities.
Thus every starred stage is an exact prefix of $A$, and $A$ is an exact
prefix of both $Y$ and $T$.  The PE ordered-union equation in
\noderef{HierarchyBase} and \noderef{OrderedTreeUnion} also puts $T_n$
and $R$ inside $A$.  Restricting
\noderef{ColoredVertexSetElementary} from $Y$ makes $A$ and every starred
stage elementary.

For any color missing on an exact prefix $U$ of $T$, the finite set of
positions in $T$ missing that color is nonempty.  Its least position is
at most the position of a witness in $U$, so its vertex lies in $U$ and
really misses the color.  This uses the defining minimum in
\noderef{HierarchyFirstMissingPosition}, not its fallback value.
Let $b$ be the resulting first witness for $\beta$.  The case assumption
makes $\alpha$ missing on $A$, so $a\in A$.  If $\beta$ is missing on $A$,
apply \noderef{ChangeLemmaBothInside}.  Its conditional ordering premise
is vacuous since both colors are missing on $A$; case 1 supplies its
positive-divider premise, and case 2 supplies either missingness on
$T_n$ or disjointness of $\{\alpha,\beta\}$ from $D_n$.
It follows that the components through $a,b$ agree.  Hence $a$ is outside
$Q$, and $b$ is also outside $Q$ whenever $\beta$ is missing on $A$.

Elementarity now gives missing-set preservation on $A$.  Outside $Q$ it
is exactly \noderef{TwoColorComponentSwapMissingBehavior}.  At a vertex
$v\in A$ lying in $Q$, $\alpha$ cannot be missing: uniqueness of its
missing vertex would give $v=a$.  Nor can $\beta$ be missing: a witness
there would make it missing on $A$ and give $v=b$.  Both colors are
therefore present at $v$.  Each of their incident edge nodes belongs to
$Q$, by the adjacency relation in \noderef{TwoColorIncidenceGraph}, and
the swap exchanges their colors.  They both remain present.  All other
missing-color memberships are unchanged by
\noderef{TwoColorComponentSwapMissingBehavior}.  This proves equality of
the entire missing sets.  Properness is
\noderef{TwoColorComponentSwapProper}; the definition
\noderef{TwoColorComponentSwap} fixes every uncolored edge and never
changes a colored value to an uncolored one, so
\noderef{IsOneEdgeUncolored} persists.

We establish interchangeability at each starred stage carrying a missing
closed color.  At a positive stage use \noderef{step1}, with the context's
lower package, fewer-divider induction and large bound.  At stage zero
in PE use clause (iii) of \noderef{interchange}, with reflexive
\noderef{CombinedStableColoring}.  At stage zero in non-PE the core is
$T_n$, by \noderef{HierarchyBase}.  Inspect the actual terminal
\noderef{TashkinovSeriesStep}.  SE supplies
\noderef{MaximumDefectiveChoice} and its selected stable coloring;
the current-color equation identifies that coloring with $\varphi$.
RE identifies $\varphi$ with the previous coloring, so the needed
\noderef{StableColoring} is reflexive.  If $n=1$ use the
interchangeability conclusion of \noderef{TechnicalTheoremBase} on the
initial closure.  If $n>1$ use \noderef{CertifiedETTRestriction} at rung
$n-1$ and lower assertion (ii): its exact tree, computed stored set, and
coloring agreements transport the SE or RE stability premise just obtained.
In each instance the first-slot color is genuinely missing.

A component rooted at a vertex missing one of the two colors is a path
component, including a singleton: \noderef{MissingColorIncidenceEndpoint}
bounds the root's degree by one.  If $Q$ met an interchangeable tree
containing an outside missing witness, the defining uniqueness in
\noderef{TreeInterchangeable} would identify the two path components,
a contradiction.  Applying this observation to the first witness for
$c$ proves the final outside-recoloring conclusion.  For $c=\beta$ the
witness is outside by the component equality already proved.  Symmetry of
\noderef{TwoColorComponent}, \noderef{TwoColorPathComponent}, and
\noderef{TreeInterchangeable} handles the reversed color order.
The singleton root set $\{z\}$ then satisfies the complete definition of
\noderef{OutsideKempeRecoloring}.

Here is the descent needed for base stability.  Suppose an affected color
$c$ missing on $A$ is closed at some starred stage, and choose its largest
closed index $r$.  Choose the least starred index $s$ at which $c$ is
missing.  If $s>r$, then $s>0$ and its missing witness is new relative to
the preceding starred stage.  By
\noderef{GoodHierarchyNewMissingClosed} and elementarity, $c$ is closed
at $s$, contradicting the choice of $r$.  Thus $s\le r$, and $c$ is
missing at $r$.  The preceding outside conclusion puts $Q$ outside
$A_r$, hence outside the core.  Apply
\noderef{OutsideKempeRecoloringStable} to $T_n$ with control $D_n$,
and in PE to $R$ with empty control, using the vertex-set inclusions in
the core.  These are exactly the asserted base stability relations.

It remains that no affected color missing on $A$ is closed on any starred
stage.  If $q>0$, apply the no-closed conclusion of \noderef{rutcor}
to $\alpha$.  Its predecessor-elementarity premise follows from
\noderef{TechnicalTheoremBase} at $n=1$, or from the maximum predecessor
returned by \noderef{CertifiedETTPredecessor} and lower assertion (i)
at $n>1$.  Its other elementarity premise is precisely the assumed
elementarity of $Y$.  This conclusion gives PE and
$a\in T_n-R$.  If $q=0$, a non-PE terminal constructor would instead
make every missing color on $A=T_n$ closed: RE uses previous closedness
and its unchanged coloring; SE transports previous closedness through its
selected stable coloring.  For the latter transport, the reverse
incident-edge clause of \noderef{sc1} fixes any hypothetical new boundary
edge whose color is tree-missing; the missing-set clause gives an old
missing witness and contradicts previous closedness.  Since $\alpha$
is missing on $A$, this rules out non-PE.

We are therefore in PE.  The first witness $a$ belongs to the core, either
by $q=0$ or by the no-closed conclusion just obtained.  It cannot belong
to $T_n\cap R$: \noderef{ParallelCommonMissingClosed}, with the same
predecessor elementarity, would make $\alpha$ closed on the core.
The ordered-union vertex set makes the following alternatives exhaustive.
If $a\in T_n$, clause (ii) of \noderef{interchange} makes $T_n$
interchangeable, so $Q$ is outside $T_n$.  If $\beta$ is missing on $R$,
clause (iv) makes the entire core interchangeable and gives the result
by outside stability.  Otherwise neither affected color is missing on
$R$: for $\alpha$ this follows from uniqueness of its witness and
$a\notin R$, and for $\beta$ it is the current alternative.
All edges whose colors are missing on $R$ are thus unaffected.
Together with the already established missing-set equality, this gives
$(R,\varnothing,\varphi)$-stability, while outside stability gives
$(T_n,D_n,\varphi)$-stability.

If $a\notin T_n$, it lies in $R$.  Clause (i) of
\noderef{interchange} makes $R$ interchangeable and puts $Q$ outside it.
If $\beta$ is missing on $T_n$, use the reversed clause (iv) for outside
stability on the entire core.  Otherwise neither affected color is
missing on $T_n$.  Here $q$ must be zero, since the positive-divider
no-closed conclusion would put $a$ in $T_n$.  The case-2 alternative
$\alpha\in\overline\varphi(T_n)$ is impossible, so its other alternative
gives $\alpha,\beta\notin D_n$.  Every controlled incident-edge color in
$\overline\varphi(T_n)\cup D_n$ is therefore unaffected; missing sets were
already proved equal.  This is $(T_n,D_n,\varphi)$-stability.
Combine it with outside stability on $R$.  All branches now give the
claimed base relation, without asserting that $Q$ avoids all of $A$.
\end{proof}

